\section{A11yPreserve Method}
\label{sec:method}

A11yPreserve treats accessibility preservation as a source-to-destination comparison problem. The method begins with a document whose target accessibility properties are known, runs that exact source through a named conversion pipeline, extracts accessibility-relevant destination structure, and classifies the correspondence. This framing is intentionally narrower than asking whether a PDF is accessible in every respect. A destination-only checker can report tags or metadata that are useful in their own right, but it cannot by itself establish whether an author-provided source property was retained.

\subsection{Source fixtures and ground truth}

The frozen corpus contains 13 small DOCX fixtures. The fixtures cover heading hierarchy, image alternative text, list semantics, table/header structure, document language, inline language changes, decorative image state, hyperlinks, document title and metadata, multi-level table headers, footnote association, equation/math semantics, and figure/caption association. They were created by controlled generation scripts using \texttt{python-docx} together with explicit WordprocessingML constructs where required. Each fixture isolates a target property and avoids unnecessary layout complexity.

Each fixture has a machine-readable preservation contract under \texttt{contracts/}. The frozen source-grounding chain used three artifacts: the source manifest, raw OOXML inspection, and the source extractor's accessibility intermediate representation (ASIR). The strengthening pass added a fourth, separate audit path: an independent source-oracle implementation performs fresh ZIP/XML assertions without importing ASIR, fixture-generation, manifest-generation, or comparator code; its per-fixture certificates are retained under \texttt{evidence/source/}. OOXML checks inspect the relevant document parts, including document XML, run properties, numbering, relationships, and the footnotes part where applicable. The ASIR remains an implementation-grounded oracle, not an independent implementation or independent human expert validation: it is part of the same analysis toolchain. Agreement among the frozen manifest, raw OOXML, independent oracle, and ASIR provides a stronger cross-check against extractor-only errors, but does not remove the limitation that no external expert panel validated every contract. This establishes encoded source ground truth; it does not claim validation by disabled users. The corpus manifest records file size, SHA-256, feature family, and verification status.
The independent OOXML oracle agreed with all 13 frozen source contracts.

\subsection{Conversion pipelines}

The first condition used LibreOffice 26.2.6.3, build \texttt{8221e31b3ac356a1623c672912a3d2b492f7e3d1}, on Windows \texttt{10.0.26200.0}. Conversion used native headless \texttt{soffice.com writer\_pdf\_Export} with \texttt{UseTaggedPDF=true}, \texttt{PDFUACompliance=true}, and \texttt{SelectPdfVersion=1}. The second condition used an authenticated Google Docs web workflow on the same Windows environment with Chrome 153.0.8010.53: upload the exact frozen DOCX and choose \textit{File} $\rightarrow$ \textit{Download} $\rightarrow$ \textit{PDF Document (.pdf)}. The observed Google PDF producer was \texttt{Skia/PDF m156}. No print-to-PDF path, screenshot, manual content repair, or post-export accessibility repair was used. The official LibreOffice documentation describes its PDF/UA export option \cite{libreoffice-pdfua}; Google's download documentation describes the selected-file-type download workflow \cite{google-docs-download}.

The Google condition is described throughout as the \emph{Google Docs DOCX-to-PDF conversion pipeline}. It includes DOCX import, Google's internal document representation, and PDF generation. The backend version cannot be fully pinned, so the experiment date matters and a rerun may change after a service update; the workflow, browser, operating system, observed producer, and environment were recorded. The experiment does not isolate which internal stage caused an observed change. Microsoft Word was not included in the denominator because it was not a completed primary pipeline.

\subsection{Destination extraction and validation}

All 26 outputs were checked for existence, non-zero size, parseability, expected page count, structural tagging, and a PDF structure tree. The extractor inspected PDF structure elements and roles relevant to the fixture contracts: headings, figures and \texttt{/Alt}, list containers and \texttt{/ListNumbering}, tables and header cells, \texttt{/Lang}, link annotations, metadata, \texttt{/Note}, \texttt{/Formula}, and \texttt{/Caption}. Evidence was retained as structured JSON and raw object observations. Poppler \texttt{pdftotext} was used as an independent secondary check when marked-content text was not recoverable from the structural parser.

The extractor is deliberately not treated as infallible. For example, the LibreOffice F11 output contains a \texttt{/Note} structure, but the available extraction evidence does not establish the note-body association. That case is retained as \measured{} rather than converted into a positive or negative preservation claim. For F06, F11, and the difficult partial cases, an independent triangulation path uses PyMuPDF page/text/link extraction and serialized-object inspection in addition to pypdf structure traversal. This reduces, but does not eliminate, parser uncertainty.

\subsection{Differential comparison}

The \texttt{a11ydiff} implementation loads the source ASIR and destination extraction for each fixture/pipeline pair. It records the source value, destination value, destination representation, evidence paths, source/output hashes, confidence, and a classification. The feature contracts and classification rules were fixed before the final conversion results were reviewed. A separate synthetic calibration set exercises positive, negative, and alternate-equivalent destination controls; its diagnostic counts are not population estimates and reveal conservative false-negative/indeterminate behavior for some alternate representations. The implementation also records a secondary destination-accessibility axis. The final evidence-aware result labels are defined below.

\begin{description}
\item[\verifiedpreserved] The destination expresses the contracted property exactly or through a verified equivalent representation, with independent source and destination evidence.
\item[\observedpartial] Some contract components survive and at least one required representable component is directly demonstrated to differ or be absent.
\item[\altered] The feature remains present, but an author-relevant value changes, such as alt-text wording, locale specificity, or title metadata.
\item[\confirmedlost] The source property is valid and representable in the destination, but the required destination representation is absent and independent checks support that absence.
\item[\unresolved] Relevant destination structure exists, but available evidence cannot establish whether the source contract survived through an alternate PDF representation; absence has not been independently demonstrated.
\item[\measured] The output is valid and relevant structure is present, but the available evidence cannot resolve the contract reliably.
\end{description}

The historical V1 result set used five categories: PRESERVED, PARTIALLY\_PRESERVED, ALTERED, LOST, and MEASUREMENT\_ERROR. The primary manuscript result model is the six-way evidence-aware taxonomy above; V1 remains archived for provenance. The source and destination files did not change.

\subsection{Independent verification and freezing}

Non-preserved cases were reviewed against the manifest, raw OOXML, independent source certificates, source ASIR, output validation, primary PDF extraction, full structure traversal, and independent parser/object checks where appropriate. The two \confirmedlost{} cases were retained only after two destination evidence paths corroborated absence of the required machine-verifiable representation; visible text survival is not treated as preservation of the associated structure. Source fixtures and PDF outputs were SHA-256 frozen. Phase 6 independently recomputed all 13 source hashes and 26 output hashes, normalized the primary and \texttt{a11ydiff} classification records, and reconstructed all 26 rows without discrepancy. The V2 evidence-aware result file is derived from the frozen V1 records and does not change fixtures, PDFs, manifests, hashes, or denominators.


